\documentclass[12pt]{article}

\usepackage[utf8]{inputenc}
\usepackage{amsmath, amssymb, amsthm}
\usepackage{geometry}
\geometry{a4paper, margin=2.5cm}

\title{Teorema de Agregare pentru Familii de Mulțimi Independente}
\author{}
\date{}

\theoremstyle{plain}
\newtheorem{theorem}{Teoremă}

\begin{document}

\maketitle

\section{Definiție: Familii de mulțimi independente}

Fie $(\Omega, \mathcal{F}, \mathbb{P})$ un spațiu de probabilitate.
Familiile de mulțimi


\[
\mathcal{A}_1, \mathcal{A}_2, \dots, \mathcal{A}_n \subseteq \mathcal{F}
\]


se numesc \textbf{independente} dacă pentru orice alegere de mulțimi


\[
A_i \in \mathcal{A}_i, \qquad i = 1,2,\dots,n,
\]


avem relația:


\[
\mathbb{P}(A_1 \cap A_2 \cap \dots \cap A_n)
= \mathbb{P}(A_1)\,\mathbb{P}(A_2)\cdots\mathbb{P}(A_n).
\]



Aceasta înseamnă că orice selecție de evenimente, câte unul din fiecare familie,
formează un set de evenimente independente.

\section{Teorema de Agregare}

\begin{theorem}[Teorema de Agregare]
Fie $\mathcal{A}_1, \mathcal{A}_2, \dots, \mathcal{A}_n$ familii de mulțimi
din $\mathcal{F}$.

Presupunem că familiile sunt independente în sensul definiției de mai sus.

Atunci pentru orice alegere de mulțimi


\[
A_i \in \mathcal{A}_i, \qquad i = 1,2,\dots,n,
\]


mulțimile $A_1, A_2, \dots, A_n$ sunt independente, adică:


\[
\mathbb{P}(A_1 \cap A_2 \cap \dots \cap A_n)
= \mathbb{P}(A_1)\,\mathbb{P}(A_2)\cdots\mathbb{P}(A_n).
\]



Cu alte cuvinte, independența familiilor implică independența oricărei
selecții de evenimente câte unul din fiecare familie.
\end{theorem}

\section{Interpretare}

Teorema afirmă că:

\begin{itemize}
    \item dacă familiile de evenimente sunt independente între ele,
    \item atunci orice alegere de evenimente, câte unul din fiecare familie,
          este independentă,
    \item deci independența se ,,agregă'' de la familii la evenimente individuale.
\end{itemize}

Aceasta este o teoremă fundamentală în teoria probabilităților, deoarece
permite construirea independenței pentru structuri complexe pornind de la
familii independente.

\section{Exemplu}

Dacă $X$ și $Y$ sunt variabile aleatoare independente, atunci familiile:


\[
\mathcal{A}_X = \{ X \in B : B \in \mathcal{B}(\mathbb{R}) \}, \qquad
\mathcal{A}_Y = \{ Y \in C : C \in \mathcal{B}(\mathbb{R}) \}
\]


sunt independente.

Prin teorema de agregare, pentru orice mulțimi măsurabile $B$ și $C$,
evenimentele


\[
\{X \in B\}, \qquad \{Y \in C\}
\]


sunt independente.

\end{document}
